\subsection[The case when $p=2r_1r_2$]{The case when $\boldsymbol{p=2r_1r_2}$}

Assume that $p=2r_1r_2$ with $r_1$, $r_2$ distinct odd primes. In this case, there exists an unique integer $x\in \{ 1, \ldots, r_1r_2-2\}$ such that $\muu{x}=1$. Then $x$ is even and verifies either
\begin{gather*}\begin{cases}
x\equiv -2 &\pmod{r_1}, \\
x\equiv 0 & \pmod{r_2},
\end{cases} \qquad \text{or}\qquad \begin{cases}
x\equiv 0 &\pmod{r_1}, \\
x\equiv -2 & \pmod{r_2}.
\end{cases}\end{gather*}
We begin by stating a technical lemma whose proof will be the subject of the next subsection:
\begin{Lemma}\label{GA}
 If $(x,x,x)$ is $p$-admissible and $r_1,r_2>37$, then we have the following inequality
\begin{gather*} \sixj{x}{x}{2}{x}{x}{0} \sixj{x}{x}{4}{x}{x}{x} \neq \sixj{x}{x}{4}{x}{x}{0}\sixj{x}{x}{2}{x}{x}{x}.\end{gather*}
\end{Lemma}

 Consider the two banded graphs $\Gamma_1= \Thetagraph$ and $\Gamma_2=\Eyesgraph$. Denote by $b_1, b_2 \in V_p(\Sigma_2)$ the two vectors $b_1:= \eyesgraph{$x$}{0}{0}$ and $b_2 := \eyesgraph{0}{0}{$x$}$.

\begin{Lemma}\label{lemma_2r1r2}
Suppose $p=2r_1r_2$, with $r_1$, $r_2$ distinct odd primes, such that either $(x,x,x)$ is not $p$-admissible, or $r_1,r_2>37$, then
\begin{gather*} W_{[0]}(\Gamma_1) \bigcap W_{[0]}(\Gamma_2) \subset \Span ( v_0, b_1, b_2 ).\end{gather*}
\end{Lemma}

\begin{proof}The subspaces $W_{[0]}(\Gamma)$ are spanned by the vectors associated to colorings of the edges of~$\Gamma$ by the elements $0$ and $x$. If $(x,x,x)$ is not $p$-admissible, then these colorings of $\Gamma_2$ correspond to the elements $v_0$, $b_1$, $b_2$ and the proof is immediate. Suppose $(x,x,x)$ is $p$-admissible. We need to show that if $v=\alpha \thetagraph{$x$}{0}{$x$} + \beta \thetagraph{$x$}{$x$}{$x$} \in W_{[0]}(\Gamma_2)$, then $\alpha=\beta = 0$.

Lemma \ref{fusion_rule} implies that
\begin{gather*}
v= \left( \alpha \sixj{x}{x}{2}{x}{x}{0} +\beta \sixj{x}{x}{2}{x}{x}{x} \right) \eyesgraph{x}{2}{x}\\
\hphantom{v=}{}+ \left( \alpha \sixj{x}{x}{4}{x}{x}{0} +\beta \sixj{x}{x}{4}{x}{x}{x} \right) \eyesgraph{x}{4}{x} + v',
\end{gather*}
where $v'$ is a linear combination of vectors of the form $\eyesgraph{$x$}{$i$}{$x$}$ for $i\neq 2,4$. Since $v\in W_{[0]}(\Gamma_2)$, the two coefficients in front of $\eyesgraph{$x$}{2}{$x$}$ and $\eyesgraph{$x$}{4}{$x$}$ vanish. Hence we get
\begin{gather*}\begin{pmatrix} \sixj{x}{x}{2}{x}{x}{0} & \sixj{x}{x}{2}{x}{x}{x} \\ \sixj{x}{x}{4}{x}{x}{0} & \sixj{x}{x}{4}{x}{x}{x} \end{pmatrix} \cdot \begin{pmatrix} \alpha \\ \beta \end{pmatrix} = \begin{pmatrix} 0\\0 \end{pmatrix}.\end{gather*}

Lemma \ref{GA} concludes the proof.
\end{proof}

\begin{Lemma}\label{lemmalacon}
There exists an element $a \in \mathcal{A}_{2r_1r_2,1}$ such that
\begin{gather*}
a\cdot u_0 = u_x, \qquad a \cdot u_x = u_0.
\end{gather*}
\end{Lemma}

\begin{proof} It is enough to show that there exists an operator $\psi \in (\mathcal{A}_{2r_1r_2,1})'$ such that
\begin{gather*} \psi \cdot u_0 = u_x \qquad \text{and} \qquad \psi \cdot u_x =u_0.\end{gather*}

The cyclicity of $u_0$, provided by Lemma \ref{lemma_genusone1}, implies the existence of $a\in\mathcal{A}_{2r_1r_2,1}$ such that
$a\cdot u_0 = u_x$. If such a $\psi$ exists, then
\begin{gather*} a\cdot u_x = a \circ \psi \cdot u_0 = \psi \circ a \cdot u_0 = u_0.\end{gather*}

The operator $\psi$ is defined as follows. Given $i\in \{0, \ldots, r_1r_2-2\}$, only one of the following two cases occurs:
\begin{enumerate}\itemsep=0pt
\item Either there exists $j\in\{0, \ldots, r_1r_2-2 \}$ so that
\begin{gather*} \begin{cases}
j\equiv i & \pmod{2r_1}, \\
j\equiv -i-2 & \pmod{r_2},
\end{cases}\end{gather*}
 and we set $\psi(u_i):=u_j$.
\item Or there exists $j\in\{0, \ldots, r_1r_2-2 \}$ so that
\begin{gather*} \begin{cases}
j\equiv i & \pmod{2r_2}, \\
j\equiv -i-2& \pmod{r_1},
\end{cases}\end{gather*}
 and we set $\psi(u_i):=-u_j$.
\end{enumerate}
A straightforward computation shows that $\psi$ commutes with $\rho_{p,1}(T)$ and $\rho_{p,1}(S)$ and either $\psi$ or $-\psi$ sends $u_0$ to $u_x$.
\end{proof}

We define two operators $a\otimes \mathds{1}, \mathds{1}\otimes a \in \mathcal{A}_{p, 2}$ as follows. Consider $\Gamma_2, \Gamma'_2 \subset S^3$ the two entangled genus $2$ banded graphs of Fig.~\ref{graph_embedding}. Denote by $\Gamma_1\subset \Gamma_2$ and $\Gamma'_1\subset \Gamma'_2$ the genus one subgraphs of the left picture of Fig.~\ref{graph_embedding}. Fix $H_2, H'_2$ some tubular neighborhood of $\Gamma_2$ and $\Gamma'_2$ respectively which do not intersect and $H_1\subset H_2$, $H'_1\subset H'_2$ some tubular neighborhoods of $\Gamma_1$ and $\Gamma'_1$ respectively. Identify the closure of $S^3\setminus (H_i\cup H'_i)$ with $\Sigma_i\times [0,1]$ and $V_p(\Sigma_i)$ with the space of linear combinations of framed links in $H_i$ quotiented by the kernel of the Hopf pairing induced by the Heegaard splitting $S^3= H_i \cup \overline{(S^3\setminus H_i)}$. From Lemma~\ref{lemmalacon}, there exists some $a\in \mathcal{A}_{p,1}$ such that $a\cdot u_x = u_0$ and $a\cdot u_0=u_x$. Let $\mathcal{L}$ be a linear combination of aligned links in $\overline{S^3\setminus (H_1\cup H'_1)}\cong \Sigma_1\times [0,1]$ such that $\Add_{p,1}\left(\mathcal{L}(\omega)\right)=a$. Composing eventually by an isotopy, we can suppose that the framed parallel links of $\mathcal{L}$ are in $\overline{S^3\setminus (H_2\cup H'_2)}\cong \Sigma_2\times [0,1]$. We denote by $a\otimes \mathds{1}:=\Add_{p,2}\left(\mathcal{L}(\omega)\right) \in \mathcal{A}_{p,2}$ the resulting operator. Note that this definition is not canonical since it depends on the choice of $a$ and on how the middle edges of~$\Gamma_2$ and~$\Gamma'_2$ are entangled with the framed links of $\mathcal{L}$. Nevertheless, by definition this operator satisfies $(a\otimes \mathds{1})\cdot v_0 = b_1$ and $(a\otimes \mathds{1})^2 \cdot v_0 = v_0$. Similarly, by considering the genus one subgraphs of $\Gamma_2$ and $\Gamma'_2$ of the right picture of Fig.~\ref{graph_embedding}, we define an operator $\mathds{1}\otimes a \in \mathcal{A}_{p,2}$ such that $(\mathds{1}\otimes a)\cdot v_0 = b_2$ and $(\mathds{1}\otimes a)^2 \cdot v_0 = v_0$. We further suppose that the aligned framed links defining these two operators are not entangled so the two operators $a\otimes \mathds{1}$ and $\mathds{1}\otimes a$ commute. We eventually define $a\otimes a := (a\otimes \mathds{1})\circ (\mathds{1}\otimes a)$ which satisfies $(a\otimes a)\cdot v_0 = \eyesgraph{$x$}{0}{$x$}$ and $(a \otimes a)^2\cdot v_0 = v_0$.

 \begin{figure}[!h]\centering
\includegraphics[width=12cm]{Korinman-Fig6}
\caption{The genus $2$ graphs $\Gamma_2$ and $\Gamma'_2$ and the genus one subgraphs $\Gamma_1$ and $\Gamma'_1$ used to define $a\otimes \mathds{1}$ on the left and $\mathds{1}\otimes a$ on the right.}\label{graph_embedding}
\end{figure}

\begin{proof}[Proof of Theorem \ref{main_th} when $p=2r_1r_2$]
Recall we defined a linear map $f\colon (\mathcal{A}_{p,2})' \hookrightarrow V_p(\Sigma_2)$ by $f(\theta)=\theta\cdot v_0$ and that the image $f\left( (\mathcal{A}_{p,2})' \right) \subset W_{[0]}(\Gamma_2)$ is invariant under the orientation-preserving homeomorphisms of the handlebody $H_2$. Using Lemma~\ref{lemma_2r1r2} and the fact that the vectors in the image of $f$ must be invariant under a homeomorphism permuting the two handles of $H_2$, we have the inclusion
\begin{gather*} f ( (\mathcal{A}_{p,2})' ) \subset \Span ( v_0, b_1+b_2 ).\end{gather*}

We need to show that $f( (\mathcal{A}_{p,2})' )= \mathbb{C}\cdot v_0$. By contradiction, suppose that there exists $\theta \in (\mathcal{A}_{p,2})'$ such that $\theta \cdot v_0 = b_1+b_2= (a\otimes \mathds{1} + \mathds{1}\otimes a) \cdot v_0$. Then we have
\begin{gather*}
\theta ^2 \cdot v_0 =( a \otimes \mathds{1} + \mathds{1}\otimes a )^2\cdot v_0= 2 v_0 + 2\eyesgraph{$x$}{0}{$x$}.
\end{gather*}
Thus $\theta^2\cdot v_0$ does not belong to $ \Span ( v_0, b_1+b_2 ) $. This contradicts the fact that $\theta^2 \in (\mathcal{A}_{p,2})'$.
\end{proof}

\subsection{Proof of Lemma \ref{GA}}

In this subsection we consider $p=2r_1r_2$, with $r_1, r_2>37$ two distinct odd primes. We suppose that there exists $x\in\{1,\ldots, r_1r_2-2\}$ so that $(x,x,x)$ is $p$-admissible and
\begin{gather*}\begin{cases}
x\equiv 0 & \pmod{2r_1}, \\
x\equiv -2 & \pmod{r_2}.
\end{cases}
\end{gather*}

We also choose some primitive $r_1$-{th} and $r_2$-{th} roots of unity $A_1$ and $A_2$, such that $A^2=A_1A_2$. In particular, we have $A^{2x}=A_2^{-2}$.

The goal of this subsection is to show that
\begin{gather*} D:= \sixj{x}{x}{2}{x}{x}{0}\sixj{x}{x}{4}{x}{x}{x} - \sixj{x}{x}{4}{x}{x}{0}\sixj{x}{x}{2}{x}{x}{x} \neq 0.\end{gather*}

A straightforward computation, using the formula of the recoupling coefficients~\cite{MV}, gives
\begin{gather*}
D = (-1)^{\frac{x}{2}+1}\frac{[3][5]![x] \left[\frac{3}{2}x+1 \right]!(\left[\frac{x}{2}\right]!)^3}{[2][x+3]!([x+2]!)^2[x+3] \left[\frac{x}{2}+1\right]}
\left( \left[\frac{x}{2}-1\right]^2 \left[\frac{x}{2}\right]^2\left[\frac{x}{2}+1\right] \right.\\
\hphantom{D =}{} - [2]^2\left[\frac{x}{2}\right]^3\left[\frac{3}{2}+2\right] \left[\frac{x}{2}+1\right] + \left[\frac{x}{2}-1\right]\left[\frac{x}{2}\right] \left[\frac{x}{2}+1\right] \left[\frac{x}{2}+2\right] \left[\frac{x}{2}+3\right]\\
\left.\hphantom{D =}{} + [x-1][x+3] \left[\frac{x}{2}\right]^2\left[\frac{x}{2}+1\right] - [x-1]\left[\frac{3}{2}x+2\right][x+3] \right) \\
\hphantom{D}{} = (-1)^{\frac{x}{2}+1}\frac{[3][5]![x]\left[\frac{3}{2}x+1\right]!(\left[\frac{x}{2}\right]!)^3}{[2][x+3]!([x+2]!)^2[x+3]\left[\frac{x}{2}+1\right](A^2-A^{-2})^7A_1^{10}A_2^{9}}\cdot P(A_1,A_2),
\end{gather*}
where $P$ is defined by
\begin{gather*}
P(x,y):= x^{20}y^{16}-x^{17}y^{17}-x^{16}y^{18}+x^{19}y^{13}-4x^{18}y^{14}+3x^{17}y^{15}+2x^{15}y^{17}-x^{19}y^{11}\\
\hphantom{P(x,y):=}{} -5x^{17}y^{13}-4x^{15}y^{15}+4x^{14}y^{16}-2x^{13}y^{17}-x^{18}y^{10} +2x^{17}y^{11}+6x^{16}y^{12} \\
\hphantom{P(x,y):=}{}+2x^{15}y^{13}-x^{14}y^{14}+2x^{13}y^{15}+x^{11}y^{17}+2x^{15}y^{11} +x^{14}y^{12}+x^{13}y^{13}-6x^{12}y^{14} \\
\hphantom{P(x,y):=}{}+x^{10}y^{16}+x^{17}y^7+4x^{16}y^8-x^{15}y^9-4x^{14}y^{10}+4x^{12}y^{12}+x^{15}y^7+x^{13}y^9 \\
\hphantom{P(x,y):=}{}-4x^{12}y^{10}+x^{11}y^{11}+4x^{10}y^{12}-4x^8y^{14}-x^7y^{15}-2x^{15}y^5-6x^{14}y^6-4x^{13}y^7 \\
\hphantom{P(x,y):=}{}+2x^{12}y^8 -8x^{11}y^9-6x^{10}y^{10}-6x^9y^{11}-x^7y^{13} +x^{13}y^5+6x^{11}y^7+6x^{10}y^8 \\
\hphantom{P(x,y):=}{}+8x^9y^9-2x^8y^{10}+4x^7y^{11}+6x^6y^{12}+2x^5y^{13}+x^{13}y^3+4x^{12}y^4 -4x^{10}y^6 \\
\hphantom{P(x,y):=}{}-x^9y^7+4x^8y^8 -x^7y^9-x^5y^{11}-4x^8y^6+4x^6y^8+x^5y^9-4x^4y^{10}-x^3y^{11} \\
\hphantom{P(x,y):=}{}-x^{10}y^2+6x^8y^4-x^7y^5-x^6y^6-2x^5y^7-x^9y-2x^7y^3 +x^6y^4-2x^5y^5\\
\hphantom{P(x,y):=}{}-6x^4y^6-2x^3y^7+x^2y^8+2x^7y-4x^6y^2+4x^5y^3+5x^3y^5+xy^7-2x^5y\\
\hphantom{P(x,y):=}{}-3x^3y^3+4x^2y^4-xy^5+x^4+x^3y-y^2.
\end{gather*}

Note that $P$ does not depend neither on $r_1$, $r_2$ nor on $x$. We have to show that $P(A_1,A_2)\neq 0$. By contradiction, suppose that $P(A_1, A_2)=0$ and define $H(X):=P(X^{r_2}, X^{r_1})\in \mathbb{Z}[X]$. Since~$P$ has integral coefficients the formula $P(A_1,A_2)=0$ holds for any roots of unity $A_1$ and $A_2$ of order $r_1$ and $r_2$ respectively. In particular, we have $H(A^2)=0$ and $H(X)$ is divisible by the product $\phi_{r_1}(X)\phi_{r_2}(X)$ of the two cyclotomic polynomials of order $r_1$ and $r_2$. Thus its degree satisfies $\deg (H) \geq (r_1-1)(r_2-1)$. Using the above formula for $P$ we deduce that
\begin{gather*}
 (r_1-1)(r_2-1) \leq \deg (H) = \begin{cases}
20r_2+16r_1, & \text{if} \ 2r_2\geq r_1, \\
16r_2 +18r_1, &\text{if} \ 2r_2\leq r_1,
\end{cases}\\
\qquad {} \Leftrightarrow
 \begin{cases}
(r_1-21)(r_2-17)\leq 356, & \text{if} \ 2r_2\geq r_1, \\
(r_1-17)(r_2-19)\leq 352, &\text{if} \ 2r_2\leq r_1.
\end{cases}
\end{gather*}

Since $r_1, r_2> 37$ are prime, then $r_1, r_2\geq 41$. We conclude by noticing that none of the two equations of the above system does have solutions satisfying $r_1, r_2\geq 41$.

